\paragraph{Retrieval-stability measures.}
For each record we computed, from the final retrieved document set (ten training documents): (i) label agreement, the proportion of retrieved documents whose label equals the predicted class; (ii) top-1 agreement; (iii) majority-label share and the number of distinct labels (evidence consistency); and (iv) rank stability, the change in retrieval-label entropy and in the majority label when only the top 3, 5 or 7 documents are used. Depth-prefix analyses truncate the saved evidence and do not re-run retrieval. Label agreement with the predicted class is not independent of the predictor, because the prediction is aggregated from the same documents. Error-detection AUROC was computed with 95\% group-bootstrap intervals (2000 replicates), and all measures were analysed post hoc and descriptively, with no thresholds tuned.
\paragraph{Retrieval-configuration sensitivity.}
On the development partition we changed one setting at a time (\texttt{rrf\_top\_n} 5/10/15; \texttt{topk\_retrieve} 10/20/40; \texttt{cascade\_top\_k\_classes} 3/5/7; shortlist filtering on/off; query word dropout 10\%/20\%; query truncation to 80\%), with the fixed seed, models and revisions of the base configuration. The base configuration was verified to reproduce the saved validation output before any variant was run. Intervals are group-bootstrap (1000 replicates). No setting was selected or tuned.
